
\section{Introduction}
\label{sec:introduction}

Many software engineering tasks come in pairs. Obfuscation and
deobfuscation~\cite{yadegari2015generic}, translation and
back-translation~\cite{lachaux2021unsupervised,pan2024lost}, and compilation and
decompilation~\cite{tan2024llm4decompile}each require a model to move between two related representations of the same program. Yet fine-tuning practice evaluates only the direction on which the model was trained. A given LLM can therefore be judged
successful after it improves on the trained transformation, even when it has lost
the inverse capability that developers rely on for inspection, migration,
debugging, and reverse engineering.

This paper asks whether forward fine-tuning preserves the reverse capability of LLMs. We ground the question in three paired settings: \wcircle{1} obfuscation and
deobfuscation, \wcircle{2} translation and back-translation, and \wcircle{3} compilation and
decompilation. Obfuscation provides the primary controlled testbed. Its forward
direction is cheap to supervise with existing
tools~\cite{superblaubeere27obfuscator}, whereas its reverse direction is hard to
supervise, because many deobfuscated programs are semantically acceptable and the
original lexical choices cannot be recovered
uniquely~\cite{raychev2015jsnice, bichsel2016deguard}. Translation
and compilation widen the question to two further inverse tasks that bear
directly on software modernization and reverse engineering.

The finding is negative, and ordinary forward evaluation would not expose it.
Across six LLMs, three semantics-preserving obfuscation types, and four prompting
strategies, 
standard fine-tuning raises forward accuracy to between 12.0\% and
81.3\% before self-correction, following a complexity ordering from lexical to semantic transformations, while reducing reverse (i.e., deobfuscation) accuracy up to 0\%. 
The reverse failure takes a
specific form. The fine-tuned models return semantically valid programs that stay
near-verbatim copies of the obfuscated input, so they remain at the transformed
endpoint rather than moving back toward a readable equivalent. The two additional
settings reproduce the same direction of effect. Forward-only fine-tuning on
Python-to-Java lowers Java-to-Python performance from 85.4\% to 63.3\%, and
forward-only fine-tuning on C-to-assembly lowers assembly-to-C performance from
16.5\% to 5.5\%.

The Bidirectional Reasoning Framework (BRF) makes the failure visible. It evaluates a fine-tuned model on both the
trained transformation and its inverse, then reports the gap between the two
directions. The framework does not require recovery of the exact original string.
It asks a more basic question. Once a model has been adapted to produce a forward
transformation, does it still move code in the reverse direction at all? The
behavior we observe is systematic. A fine-tuned model generates almost
exclusively in the trained direction even when the reverse prompt is explicit and
task-relevant
% \js{I am not sure to understand this sentence}\se{The goal is to say that fine-tuned models still generate in trained direction even if there are instruct not to, like deobfuscate and they still obfuscate}.

% This behavior echoes the reversal curse, which reports that an autoregressive LLM
% trained on a fact stated as ``A is B'' often fails to answer the reversed query
% ``B is A''~\cite{zhang2024directional, grosse2023influence}. The shared intuition
% is an asymmetry between the direction a model is trained on and the direction a
% user later queries. We name our phenomenon the \emph{inversion curse} to mark
% that lineage, the program-transformation counterpart of the reversal curse in
% factual recall.

This behavior echoes the reversal curse observed in natural language pretraining, which reports that an autoregressive LLM
trained on a fact stated as ``A is B'' often fails to answer the reversed query
``B is A''~\cite{zhang2024directional, grosse2023influence}. Yet, our setting differs fundamentally. We move from natural language facts to executable code. Our inverse tasks involve multiple valid solutions validated by test suites. Furthermore, while the reversal curse stems from pretraining dynamics, the inversion curse arises directly from downstream task-specific fine-tuning on software engineering tasks.

Moving the question into software engineering changes what can be asked and what
is at stake. A program transformation and its inverse are executable procedures
with preconditions, search spaces, and correctness criteria, so the inverse is no
longer a single reversed fact but a task that can be validated by compilation and
testing. Our transformations run from a lexical
identifier mapping to structural pruning and to algorithmic recovery of encrypted
literals. This progression lets us study how the asymmetry scales with task complexity.
Because the inverse tasks are executable and graded, software engineering offers a precise setting. In this setting, the phenomenon can be measured, decomposed, and connected directly to deployment.

% That setting also raises the stakes. In a deployed pipeline, a missing inverse
% transformation can silently disable deobfuscation, decompilation, rollback, or
% back-translation, even when the trained direction continues to pass its own
% tests. A concrete case is the organization that deploys a single fine-tuned code
% model as a shared assistant. When that model is adapted for one dominant task,
% such as a large translation or hardening effort, the same parameters serve every
% developer who could later rely on the inverse direction, and the inverse is rarely
% re-validated after the adaptation. The governing question for our community is
% therefore whether post-training adaptation preserves the adjacent capabilities
% that real workflows assume remain available. We answer it with direct measurement on six models and
% three transformation families, and we trace the asymmetry to ordinary downstream
% fine-tuning rather than to pretraining data order.

That setting also raises the stakes.
LLMs now perform both sides of paired tasks in practice.
Google reports LLM-assisted code migrations at industrial scale~\cite{nikolov2025google}.
Security work delegates both obfuscation and decompilation to LLMs~\cite{mohseni2024obfuscate, tan2024llm4decompile}.
Obfuscation is common in Android apps~\cite{wermke2018obfuscation, dong2018android}.
Analysts therefore need to deobfuscate them~\cite{bichsel2016deguard}.
In these workflows, the inverse also serves as a check on the forward output.
Round-trip correctness runs the same model forward and backward to validate its outputs without labels~\cite{allamanis2024roundtrip}.
A model adapted to one direction weakens this check.
Our translation results illustrate the risk.
Forward-only fine-tuning lowers Java-to-Python success from 85.4\% to 63.3\% (Section~\ref{sec:transl-compil}).
A round-trip check built on this model would therefore become less reliable.
%The failure is also silent.
%A failed inverse output still compiles and passes the tests of the source (Section~\ref{sec:reverse-eval}).
%Execution-based checks therefore raise no alarm.
Adapted checkpoints are also shared and reused with incomplete documentation~\cite{jiang2023reuse}.
A downstream user cannot see which adjacent capability the adaptation removed.
The governing question for our community is therefore whether post-training adaptation preserves the adjacent capabilities that real workflows assume.
We answer it by direct measurement on six models and three transformation families, using fine-tuning.
%We trace the asymmetry to ordinary downstream fine-tuning.


We address three research questions. First, does forward fine-tuning preserve or
improve reverse capability in the absence of reverse supervision? Second, what
evidence separates an ordinarily hard inverse task from a direction-specific
failure induced by fine-tuning? Third, can bidirectional supervision or
contrastive regularization recover reverse capability, and what does the recovery
pattern reveal about the underlying mechanism? The answers hold across model
families. Standard fine-tuning fails to preserve reverse capability. Balanced
bidirectional supervision recovers the lexical and structural inverses but not
the semantic inverse. Contrastive fine-tuning recovers reverse capability only on
the lexical transformation.

The chain-of-thought probe sharpens the mechanism. On identifier renaming, a
fine-tuned model states the correct reverse mapping in its reasoning trace and
then emits the obfuscated code unchanged. %We do not read the trace as a faithful report of internal state. 
We read it as empirical evidence that the reverse
mapping can be surfaced at inference time without being applied in the final
generation. The gap fits a generation-time directional bias acting on top of any
representational effect.

\noindent\textbf{We make the following contributions:}
\begin{itemize}[leftmargin=*,topsep=0pt]
    %{\color{red}{
    \item[\wcircle{1}] We identify the inversion curse, a directional capability loss in which forward fine-tuning of a code LLM improves the trained direction while driving reverse capability to zero.
    %}}
    \item[\wcircle{2}] We define BRF, a low-cost diagnostic that exposes the
    asymmetry by evaluating a fine-tuned model in both directions.
    \item[\wcircle{3}] We give empirical evidence that the lexical case reflects
    a reasoning-generation gap rather than absence of the reverse mapping.
    \item[\wcircle{4}] We derive implications for validating, deploying, and
    repairing fine-tuned code LLMs. We also confirm these obfuscation findings on translation and compilation tasks.
\end{itemize}

% The rest of the paper follows the argument of the result. \Cref{sec:theorical}
% defines directional capability loss and BRF.\Cref{sec:experiment_setup}
% describes the controlled testbed, datasets, models, and evaluation protocol.
% \Cref{sec:cognitive_specialization} presents the main empirical findings.
% \Cref{sec:cs_mechanism_section} studies the mechanism and the effect of
% explicit bidirectional supervision. \Cref{contrastive} uses contrastive
% fine-tuning as a probe into the representational component of the failure.
% \Cref{related_work} positions the paper against prior work, and
% \Cref{sec:discussion} discusses implications, threats to validity, and
% future directions. \Cref{sec:conclusion} concludes.
